\chapter{Geometría intrínseca}
% 23/10/26
Estudia $\displaystyle S $ desde dentro. Lo que mide Planito. Lo que sale de conocer la primera forma fundamental.
\begin{definition}[Distancia dentro de una superficie]
Dados $\displaystyle P,Q\in S $ en $\displaystyle S $ conexa se define la distancia entre $\displaystyle P $ y $\displaystyle Q $ como
\[d _{S}\left(P,Q\right)=\inf \left\{ \ell^{1}_{0}\left(\alpha \right) \; : \; \alpha: \left[0,1\right] \to S, \; \alpha\left(0\right)=P, \; \alpha\left(1\right)=Q\right\}  .\]
\end{definition}
% tabla de cualidades intrínsecas y extrínsecas
\begin{definition}[Intrínseco y extrínseco]
Una noción es \textbf{intrínseca} si se escribe como $\displaystyle \mathcal{F}\left(E, F, G\right) $. Es \textbf{extrínseca} si se escribe como $\displaystyle \mathcal{F}\left(E, F, G, e, f, g\right) $ pero no como $\displaystyle \mathcal{F}\left(E, F, G\right) $. 
\end{definition}
\section{Isometrías}
\begin{definition}[Isometría]
Una aplicación $\displaystyle F : S_{1}\to S_{2} $ diferenciable es:
\begin{enumerate}
\item \textbf{isometría local} si $\displaystyle d _{p}F : T_{p}S_{1}\to T_{F\left(p\right)}S_{2} $ es isometría lineal para todo $\displaystyle p \in S_{1} $. 
\item \textbf{isometría} si $\displaystyle F $ es isometría local y es biyectiva.
\end{enumerate}
\end{definition}
\begin{definition}[Difeomorfismo local y difeomorfismo]
Dada una aplicación diferenciable $\displaystyle F $ decimos que 
\begin{enumerate}
\item $\displaystyle F $ es difeomorfismo local si $\displaystyle d _{p}F : T_{p}S_{1}\to T_{F\left(p\right)}S_{2} $ es isomorfismo local en cada $\displaystyle p \in S_{1} $. 
\item $\displaystyle F $ es difeomorfismo si es difeomorfismo local y además es biyectiva.
\end{enumerate}
\end{definition}
\begin{observation}
\begin{enumerate}
\item A efectos de derivabilidad de objetos son equivalentes. 
	% esto es lo que se ve con la primera parte de la definición anterior
	% poner imagen
\item Si $\displaystyle F $ es difeomorfismo, entonces $\displaystyle S_{1} $ y $\displaystyle S_{2} $ son equivalentes desde el punto de vista de la geometría diferencial.
\item Si $\displaystyle F: S_{1}\to S_{2} $ es isometría local, entonces $\displaystyle d _{p}F : T_{p}S_{1} \to T_{F\left(P\right)}S_{2} $ es isometría. Es decir, $\displaystyle \|d _{p}F\left(\vec{v}\right)\|_{S_{2}}=\|\vec{v}\|_{S_{1}} $, $\displaystyle \forall \vec{v}\in T_{p}S_{1} $, $\displaystyle \forall p \in S_{1} $. 
\item Cuando se preservan las normas se preservan los ángulos, por lo que 
	\[L\left(d _{p}F\left(\vec{v}\right), d _{p}F\left(\vec{w}\right)\right)=L\left(\vec{v}, \vec{w}\right) .\]
\item Se tiene que 
	\[\ell^{b}_{a}\left(F\circ \alpha \right)=\int^{b}_{a} \|\left(F\circ \alpha \right)'\left(t\right)\| \; dt = \int^{b}_{a} \| dF\left(\alpha'\left(t\right)\right)\| \; dt = \int^{b}_{a} \|\alpha'\left(t\right)\| \; dt = \ell^{b}_{a}\left(\alpha \right) .\]
\end{enumerate}
\end{observation}

\begin{theorem}
Si $\displaystyle F $ es isometría local, existe coordenadas locales $\displaystyle \textbf{x} \left(u,v\right) $ en $\displaystyle U^{p}\subset S_{1} $, $\displaystyle \textbf{y} \left(u,v\right) $ en $\displaystyle F\left(U^{p}\right)\subset S_{2} $ (con los mismo parámetros $\displaystyle \left(u,v\right) $) tales que 
\[ \begin{pmatrix} E_{ \textbf{x} } & F _{ \textbf{x} } \\ F_{ \textbf{x} } & G _{ \textbf{x} } \end{pmatrix}=\begin{pmatrix} E_{ \textbf{y} } & F_{ \textbf{y} } \\ F_{ \textbf{y} } & G_{ \textbf{y} } \end{pmatrix} ,\]
y conversamente.
\end{theorem}

\begin{proof}
% poner foto del diagrama
% aparentemente toda la demostración es una foto
	Tenemos que existen $\displaystyle \mathcal{B}_{ \textbf{x} } = \left\{ \textbf{x} _{u}, \textbf{x} _{v}\right\}  $ y $\displaystyle \mathcal{B}_{ \textbf{y} }= \left\{ \textbf{y} _{ w}, \textbf{y} _{s}\right\}  $, de forma que 
	\[M_{ \mathcal{B}_{ \textbf{x} }}\left(I_{S_{1}}\right)=\begin{pmatrix} E_{ \textbf{x} } & F_{ \textbf{x} } \\ F_{ \textbf{x} } & G_{ \textbf{x} } \end{pmatrix} \quad \text{y} \quad M_{ \mathcal{B}_{ \textbf{y} }}\left(I_{S_{2}}\right)=\begin{pmatrix} E_{ \textbf{y} } & F_{ \textbf{y} } \\ F_{ \textbf{y} } & G_{ \textbf{y} } \end{pmatrix} .\]
	Si tomamos la base $\displaystyle \mathcal{B}_{1}= \left\{ \oplu_{p}F\left( \textbf{x} _{u}\right), d _{p}F\left( \textbf{x} _{v}\right)\right\}  $, tenemos que $\displaystyle d _{p}F $ es isometría lineal si y sólo si $\displaystyle M_{ \mathcal{B}_{ \textbf{x} }}\left(\left\langle ,  \right\rangle \right)=M_{ \mathcal{B}_{1}}\left(\left\langle ,  \right\rangle \right) $. De esta forma, 
	\[M_{ \mathcal{B}_{ \textbf{y} }}\left(\left\langle ,  \right\rangle \right) = M_{\mathcal{B}_{1}\mathcal{B}_{ \textbf{y} }}^{t}M_{\mathcal{B}_{1}}\left(\left\langle ,  \right\rangle \right)M_{\mathcal{B}_{1}\mathcal{B}_{ \textbf{y} }} .\]
Se sigue que 
\[
\begin{cases}
d _{p}F\left( \textbf{x} _{u}\right)=\frac{\partial w}{\partial u} \textbf{y} _{w} + \frac{\partial s}{\partial u} \textbf{y} _{s} \\
d _{p}F\left( \textbf{x} _{v}\right)=\frac{\partial w}{\partial v} \textbf{y} _{w} + \frac{\partial s}{\partial v} \textbf{y} _{s}
\end{cases}
.\]
Como $\displaystyle M_{ \mathcal{B}_{ \textbf{y} }\mathcal{B}_{1}} = J\tilde{F} $. Se sigue que 
\[M_{\mathcal{B}_{ \textbf{x} }}\left(\left\langle ,  \right\rangle \right)=M_{ \mathcal{B}_{1}}\left(\left\langle ,  \right\rangle \right)=\left(J\tilde{F}\right)^{t}M_{ \mathcal{B}_{ \textbf{y} }}\left(\left\langle ,  \right\rangle \right)\left(J\tilde{F}\right) .\]
Así, 
\[I_{ \textbf{x} }=\left(J\tilde{F}\right)^{t}I_{ \textbf{y} }\left(J\tilde{F}\right) .\]
Por tanto $\displaystyle I_{ \textbf{x} }= I_{ \textbf{y} } $ si tomo $\displaystyle \textbf{y} =F\circ \textbf{x}  $.
\end{proof}
\begin{observation}
\begin{enumerate}
\item Una isometría local puede no ser sobreyectiva. % considerar la foto
\item Una isometría local puede no ser inyectiva. Consideremos el plano $\displaystyle S_{1}= \left\{ \left(u,v,0\right)\right\}  $ y $\displaystyle \left\{ x^{2}+y^{2} = 1\right\}  $. Tenemos que el difeomorfismo local es $\displaystyle F\left(u,v\right)=\left(\cos u, \sin u, v\right) $. Tenemos que, cogiendo $\displaystyle \textbf{y} \left(\varphi,w\right)=\left(\cos \varphi, \sin \varphi, w\right) $, tenemos que 
	\[I_{S_{1}}=\begin{pmatrix} 1 & 0 \\ 0 & 1 \end{pmatrix}\quad \text{y} \quad I_{S_{2}}=\begin{pmatrix} r\left(v\right)^{2} & 0 \\ 0 & \|\gamma'\left(v\right)\|^{2} \end{pmatrix}=\begin{pmatrix} 1 & 0 \\ 0 & 1 \end{pmatrix} .\]
	Es isometría local. En el caso de que no sea inyectiva no se respetan las áreas. 
	%poner foto del cilindro
	% el ejemplo continua con algo de longitudes, foto de lo hecho en clase
	% no se conserva la curva más corta
\end{enumerate}
\end{observation}
\begin{prop}
Si $\displaystyle F : S_{1}\to S_{2} $ es isometría (global) entonces:
\begin{enumerate}
\item $\displaystyle Area\left(R\right)=Area\left(F\left(R\right)\right) $ para todo $\displaystyle R \subset S_{1} $ compacto. 
\item $\displaystyle d _{S_{1}}\left(p,q\right)=d _{S_{2}}\left(F\left(p\right), F\left(q\right)\right) $.
\end{enumerate}
\end{prop}
\begin{definition}
Una noción es intrínseca si es invariante por isometría local. 
\end{definition}
Estas sutilezas son de carácter global.
% 24/11/26
Sea $\displaystyle T : \R^{3}\to \R^{3} $ un movimiento rígido, esto es, $\displaystyle T\left(X\right)=AX + B $ con $\displaystyle A^{t}A = I $ y $\displaystyle A \in O_{3}\left(\R\right) $. 
\begin{definition}[Superficies congruentes]
Decimos que dos superficies $\displaystyle S_{1} $ y $\displaystyle S_{2} $ son \textbf{congruentes} si existe un movimiento rígido $\displaystyle T $ tal que $\displaystyle T\left(S_{1}\right)=S_{2} $.
\end{definition}
% poner dibujo de la foto
\begin{observation}
Si tomamos la parametrización $\displaystyle \textbf{x} \left(u,v\right) $ de $\displaystyle S_{1} $ y entonces 
\[ \textbf{y} \left(u,v\right)=T\left( \textbf{x} \left(u,v\right)\right)=A \textbf{x} \left(u,v\right)+B ,\]
es parametrización de $\displaystyle S_{2} $. En este caso podemos observar que tenemos las igualdades:
\[ \textbf{y} _{u}= A \textbf{x} _{u} , \quad \textbf{y} _{v}= A \textbf{x} _{v}, \quad \textbf{y} _{u}\times \textbf{y} _{v}= A \textbf{x} _{u} \times A \textbf{x} _{v} = \epsilon A\left( \textbf{x} _{u} \times \textbf{x} _{v}\right) .\]
Así, 
\[ N_{ \textbf{y} }= \frac{ \textbf{y} _{u} \times \textbf{y} _{v}}{\| \textbf{y} _{u} \times \textbf{y} _{v}\|}=\frac{\epsilon A\left( \textbf{x} _{u}\times \textbf{x} _{v}\right)}{\| \textbf{x} _{u}\times \textbf{x} _{v}\|}=\epsilon AN_{ \textbf{x} } .\]
También se puede ver
% hacer cálculos
\[\begin{pmatrix} E_{ \textbf{x} } & F_{ \textbf{x} } \\ F_{ \textbf{x} } & G_{ \textbf{x} } \end{pmatrix}=\begin{pmatrix} E_{ \textbf{y} } & F _{ \textbf{y} } \\ F_{ \textbf{y} } & G_{ \textbf{y} } \end{pmatrix}\quad \text{y} \quad \begin{pmatrix} e_{ \textbf{y} } & f_{ \textbf{y} } \\ f_{ \textbf{y} } & g_{ \textbf{y} } \end{pmatrix}=\pm \begin{pmatrix} e_{ \textbf{x} } & f_{ \textbf{x} } \\ f_{ \textbf{x} } & g_{ \textbf{x} } \end{pmatrix} .\]
\end{observation}
\begin{colorary}
Si $\displaystyle S_{1} $ y $\displaystyle S_{2} $ son contruentes, entonces son isométricas. 
\end{colorary}
\begin{colorary}
Las superficies congruentes comparten (con cartas adoptadas) la primera y segunda forma fundamental.
\end{colorary}
% demostraciones están en la observación anterior
% copiar el esquema de implicaciones de la foto
% que no tengan la misma primera forma fundamental no garantizan que no sean isométricas, hay que resolver el sistema $\displaystyle I_S1 \neq (Jh)^t I_S2 (Jh) $ 
% (c) es cierto para n=2 superficie, pero es un tema muy complicado
\begin{eg}[Hoja 4, Catenoide]
Es la superficie de revolución dada por la curva $\displaystyle \gamma\left(v\right)=\left(\cosh v, v\right) $. Así, podemos tomar la parametrización
\[ \textbf{x} \left(\varphi, v\right)=\left(\cosh v \cos \varphi, \cosh v\sin\varphi, v\right) .\]
Es sencillo ver que
\[I = \begin{pmatrix} \cosh^{2}v & 0 \\ 0 & \cosh^{2}v \end{pmatrix}=\begin{pmatrix} r^{2} & 0 \\ 0 & \|\gamma'\|^{2} \end{pmatrix} \quad \text{y} \quad II = \begin{pmatrix} -\frac{r}{\|\gamma'\|} & 0 \\ 0 & \frac{r''}{\|\gamma'\|} \end{pmatrix}=\begin{pmatrix} -1 & 0 \\ 0 & 1 \end{pmatrix}.\]
De esta forma, obtenemos que 
\[K = -\frac{r''}{r\|\gamma'\|^{4}} .\]
Por otro lado, 
\[\mathcal{W}=\begin{pmatrix} -\frac{1}{\cosh^{2}v} & 0 \\ 0 & \frac{1}{\cosh^{2}v} \end{pmatrix} .\]
De aquí se deduce que $\displaystyle H = 0 $ y $\displaystyle K = -\frac{1}{\cosh^{4}v} $. Se puede observar que esta superficie es minimal. 
\end{eg}
\begin{eg}[Hoja 4, Helicoide]
Se construye mediante la hélice dada por $\displaystyle \alpha\left(t\right)=\left(\cos t, \sin t, t\right) $. La podemos parametrizar por
\[ \textbf{x} \left(t, \lambda \right)=\left(\lambda \cos t, \lambda \sin t, t\right) .\]
Se puede observar que es una superficie reglada. Tenemos que 
\[ \textbf{x} _{t}=\begin{pmatrix} -\lambda\sin t \\ \lambda \cos t \\ 1 \end{pmatrix}\quad \text{y} \quad \textbf{x} _{\lambda }=\begin{pmatrix} \cos t \\ \sin t \\ 0 \end{pmatrix} .\]
Así, 
\[I = \begin{pmatrix} \lambda^{2} + 1 & 0 \\ 0 & 1\end{pmatrix} .\]
Resulta que es isométrica a la catenoide. En efecto, tomamos 
\[
\begin{cases}
t = \varphi \\
\lambda = \sinh v
\end{cases}
.\]
tenemos que, 
\[\left(t, \lambda \right)=\tilde{F}\left(\varphi, v\right)=\left(\varphi, \sinh v\right) .\]
Por tanto, $\displaystyle \textbf{y} \left(u,v\right)=\left(\sinh v \cos \varphi, \sinh v \sin \varphi, \varphi\right) $. 
% poner dibujo de la foto
Tenemos que 
\[I _{ \textbf{y} \left(\varphi, t\right)}=\left(J\tilde{F}\right)^{t}\begin{pmatrix} \lambda^{2} + 1 & 0 \\ 0 & 1 \end{pmatrix}\left(J\tilde{F}\right)=\begin{pmatrix} 1& 0 \\ 0 & \cosh v \end{pmatrix}\begin{pmatrix} \sinh^{2}v+1 & 0 \\ 0 & 1 \end{pmatrix}\begin{pmatrix} 1 & 0 \\ 0 & \cosh v \end{pmatrix}=\begin{pmatrix} \cosh^{2}v & 0 \\ 0 & \cosh^{2}v \end{pmatrix},\]
que es la primera forma fundamental del cateoride. Tenemos que $\displaystyle F : S_{1} \to S_{2} $ (con $\displaystyle S_{1} $ el cateroide y $\displaystyle S_{2}$ el helicoide) es una isometría local. Si cogemos $\displaystyle \varphi\left(0, 2\pi \right) $ y $\displaystyle v \in \left(0, \infty\right) $ tenemos que la aplicación $\displaystyle F $ es biyectiva. 
% calcular la segunda forma fundamental que viene dada en la foto
% en la foto de la pizarra hay algo con el signo mal, hacerlo desde cero y ver que da bien como en la otra foto
Las dos que acabamos de ver son isométricas, no son congruentes, comparten la misma curvatura, es decir $\displaystyle K_{S_{1}}=K_{S_{2}} $. Además, $\displaystyle H_{S_{1}}=H_{S_{2}}=0 $. La curvatura media no es intrínseca, como se verá posteriormente. Sin embargo, sí se puede ver que $\displaystyle \vec{e}_{1}=\left(1,0\right)_{\mathcal{B}_{ \textbf{x} }} $ y $\displaystyle \vec{e}_{2}=\left(0,1\right)_{\mathcal{B}_{ \textbf{x} }} $ de $\displaystyle S_{1} $. Por otro lado, $\displaystyle \vec{e}_{1}=\left(1,-1\right)_{\mathcal{B}_{ \textbf{y} }} $ y $\displaystyle \vec{e}_{2}=\left(1,1\right)_{\mathcal{B}_{ \textbf{y} }} $ en $\displaystyle S_{2} $. 
Por tanto, las direcciones principales, las direcciones asintóticas, las líneas de curvatura y las líneas asintóticas no son intrínsecas, sino extrínsecas. 
% poner esto último como teorema
\end{eg}

